\section*{Question 1}

\partquestion{(a)}A logistic regression model was fit on the \emph{LogData.xlsx} dataset. This was done
using the \emph{Logit} class from the \emph{statsmodels} library in Python.

\begin{enumerate}[label=\roman*.]
    \item After fitting the model, the following estimated values were obtained:
         \[\hat{\beta_0} = -1.9496\]
         \[\hat{\beta_1} = 0.0243\]

         Hence, the linear predictor is:
         \[\hat{\eta}(x) = -1.9496 + 0.0243x\] 
         
    \item Using the linear predictor equation provided:
         \[\log \frac{P(TT = 1 | x)}{1 - P(TT = 1 | x)} = \hat{\eta}(x),\]

         the probability that TT equals 1 given $x$ (estimated prediction probability function) is:
         \[P(TT = 1 | x) = \frac{1}{1 + e^{-\hat{\eta}(x)}} = \frac{1}{1 + e^{-(-1.9496 + 0.0242x)}}.\]

    \item Figure \ref{fig:sp_simple} illustrates the scatter plot of $TT$ against $x$:
    
          \begin{figure}[h]
            \centering
            \includegraphics[width=0.75\textwidth]{images/output_simple.png}
            \caption{Scatter Plot of $TT$ against $x$}
            \label{fig:sp_simple}
          \end{figure}

          From Figure \ref{fig:sp_simple}, the predicted probabilities of each point can be
          seen in orange with the estimated probability function plotted over the scatter points.
          The original input data points can also be seen at 0.0 and 1.0.

    \item Figure \ref{fig:cm_simple} illustrates the confusion matrix of the outcomes predicted
          by the logistic regression model:

          \begin{figure}[h]
            \centering
            \includegraphics[width=0.55\textwidth]{images/confusion_matrix_simple.png}
            \caption{Confusion Matrix}
            \label{fig:cm_simple}
          \end{figure}

          In addition to Figure \ref{fig:cm_simple}, the Somers' D value was calculated as:
          \[D_{yx} = 2 \times (AUC - 0.5) = 0.4453\]

\end{enumerate}

\partquestion{(b)}A cubic splines logistic regression model was fit on the \emph{LogData.xlsx} dataset with
knots at 50 and 80. This was done using the \emph{Logit} class from the 
\emph{statsmodels} library in Python to ensure consistency with (a).

\begin{enumerate}[label=\roman*.]
    \item The basis functions for a cubic splines logistic regression model with knots at 50 and
          80 are as follows:
          \[b_1(x) = x,\]
          \[b_2(x) = x^2,\]
          \[b_3(x) = x^3,\]
          \[b_4(x) = (x - 50)_+^3,\]
          \[b_5(x) = (x - 80)_+^3,\]

          where 
          \[(x - c)_+^3 =   \begin{cases} 
                                (x - c)^3 & x > c \\
                                0 & x \leq c.
                            \end{cases}
          \]

          The corresponding coefficient estimates from the fitted model are as follows:
          \[\hat{\beta_0} = 11.5914,\]
          \[\hat{\beta_1} = -1.0457,\]
          \[\hat{\beta_2} = 0.0276,\]
          \[\hat{\beta_3} = -0.0002,\]
          \[\hat{\beta_4} = 0.0003,\]
          \[\hat{\beta_5} = -0.0002.\]

          Hence, the predictor is as follows:
          \[\hat{\eta}(x) = 11.5914 - 1.0457x + 0.0276x^2 - 0.0002x^3 + 0.0003(x - 50)_+^3 - 0.0002(x - 80)_+^3\]

    \item Using the log-odds (logit) function:
         \[\log \frac{P(TT = 1 | x)}{1 - P(TT = 1 | x)} = \hat{\eta}(x),\]

         the probability that TT equals 1 given $x$ (estimated prediction probability function) is:
         \[P(TT = 1 | x) = \frac{1}{1 + e^{-\hat{\eta}(x)}}\]
         \[P(TT = 1 | x) = \frac{1}{1 + e^{-(11.5914 - 1.0457x + 0.0276x^2 - 0.0002x^3 + 0.0003(x - 50)_+^3 - 0.0002(x - 80)_+^3)}}.\]                      

    \newpage 

    \item Figure \ref{fig:sp_spline} illustrates the scatter plot of $TT$ against $x$:
    
          \begin{figure}[h]
            \centering
            \includegraphics[width=0.75\textwidth]{images/output_spline.png}
            \caption{Scatter Plot of $TT$ against $x$}
            \label{fig:sp_spline}
          \end{figure}

          From Figure \ref{fig:sp_spline}, the predicted probabilities of each point can be
          seen in orange with the estimated probability function plotted over the scatter points.
          The original input data points can also be seen at 0.0 and 1.0.

    \item Figure \ref{fig:cm_spline} illustrates the confusion matrix of the outcomes predicted
          by the cubic splines logistic regression model:

          \begin{figure}[h]
            \centering
            \includegraphics[width=0.55\textwidth]{images/confusion_matrix_spline.png}
            \caption{Confusion Matrix}
            \label{fig:cm_spline}
          \end{figure}

          In addition to Figure \ref{fig:cm_spline}, the Somers' D value was calculated as:
          \[D_{yx} = 2 \times (AUC - 0.5) = 0.6016\]

\end{enumerate}

\partquestion{(c)}The best model is the cubic spline logistic regression model fitted in Question 1.b.
When looking at the confusion matrix of both models (Figure \ref{fig:cm_simple} and Figure \ref{fig:cm_spline}),
it can be seen that the cubic spline logistic regression model outperforms the simple logistic regression
model producing a large reduction in False Positives from 99 to 49 and slight reduction in False Negatives
from 87 to 83. It is therefore better at classifying the training data in general.
In addition, the cubic spline logistic regression model also has a better Somers' D value
of 0.6016 against the simple logistic regression's Somers' D value of 0.4453. Somers' D measures
the model's predictive power and discrimination. Models having a Somers' D value closer to 1 are
better classifiers. Thus, the cubic spline logistic regression model is the better model. This is
expected because it is also more flexible allowing the relationship to change shape around the knots 
at 50 and 80.